\chapter{Relative Value: (Alpha Book)}

\section{Regression (\texttt{pairs/stats.py})}

A pair is modeled by OLS (via \texttt{statsmodels}) of one leg on another, encapsulated in
\texttt{RegressionResults}:

\begin{itemize}
\item Stores \texttt{intercept}, \texttt{slope\_per\_step} (the \textbf{hedge ratio}), \texttt{r2}, \texttt{n\_obs}.
\item \textbf{Residual dispersion} uses $\operatorname{std}(\text{residuals}, \mathrm{ddof}=2)$ (correct degrees of freedom for a two-parameter regression) to build \textbf{confidence bands}.
\item Signals are driven by the \textbf{residual z-score} vs. configurable entry/exit bands; mean-reversion of the residual is the exit.
\end{itemize}

\section{Spread Families (Alpha Book $\rightarrow$ Spread)}

Sector PCA spreads, spread regression, treasury/policy/local/corporate spreads, swap
spreads, bond-swap spreads, and futures term/net basis. The statistical layer
(\texttt{curves/calibration/stat.py}) supplies stationarity flags and z-scores so candidates
are only flagged when the relationship is statistically stable.

\paragraph{Review focus.}
Stationarity/cointegration testing rigor, lookback choice for the regression window,
and stability of the hedge ratio through regime shifts.

\subsection{OTR/OFR Relative Value and New-Issue Event Strategies}
\label{sec:pairs-tbondcurve-otr-ofr}

The OTR/OFR initiative is split into two strategy families that share the same
economic setup but not the same statistical contract.

Mature-pair relative value (\texttt{otr\_ofr\_rv}) remains within
\texttt{TBondCurve}/\texttt{CBondCurve} because it has a genuine pair-specific
time series suitable for z-score and MR/trend routing. The new-issue roll-pressure
strategy is represented as a unified spread type, \texttt{BondNewIssue}, with
asset distinction carried by metadata fields (for example
\texttt{asset\_class=TBond|CBond}, \texttt{issuer\_class=CGB|CDB}, and
\texttt{tenor\_bucket}).

For mature-pair RV, the trade is not an extrapolation of the fitted 1--10Y model
curve. Its economic object is the liquidity premium between current OTR and a prior
OFR issue in the same bucket.

For an OFR-minus-OTR yield spread,
\begin{equation}
	s_t = y^{\mathrm{OFR}}_t - y^{\mathrm{OTR}}_t,
\end{equation}
the executable relative-value position is long the OFR and short the OTR when the
spread is unusually wide.  The cash legs are sized DV01-neutral,
\begin{equation}
	N_{\mathrm{OTR}} = N_{\mathrm{OFR}}
	\frac{\mathrm{DV01}_{\mathrm{OFR}}}{\mathrm{DV01}_{\mathrm{OTR}}},
\end{equation}
with residual duration and convexity exposure reported separately.  This is a
sovereign liquidity/relative-value strategy, not a conventional credit-spread
strategy.  Candidate eligibility requires live two-sided quotes, adequate turnover,
short/borrow availability for the OTR leg, and an all-in expected benefit that exceeds
bid--offer, funding, borrow, and execution buffers.

\paragraph{OTR rolls and cold start policy.}
The OTR identifier changes at issuance, so a new OTR/OFR pair has no pair-specific
history on its first day.  Historical spread data remain useful, but must not be
treated as if they were observations of the new pair.  The production history is a
point-in-time, rolled pair series: each observation uses the OTR and OFR that were
eligible on that date, with the pair identifiers, issuance ages, remaining maturities,
DV01 hedge ratio, and roll date retained in the audit record.  A historical
``analogue'' distribution, conditioned on OTR issuance age and comparable 30Y
remaining-maturity buckets, may provide a conservative prior for monitoring and
entry-threshold calibration; it must not fabricate a z-score history for a new pair
or be used with future constituent knowledge.

The new pair is therefore \emph{observe-only} until it has sufficient post-roll,
executable observations.  In the baseline, it needs at least 20 valid observations
for descriptive statistics, 60 for the monthly regime classifier, and 120 for the
pair-specific MR z-score.  Before the 120-observation threshold, the candidate is
labeled \texttt{insufficient\_pair\_history}, receives no allocation, and displays its
analogue percentile, liquidity diagnostics, and accumulated live history.  A new OTR
also causes any predecessor pair to close or be explicitly rolled at the known
roll-date observation; the two histories are not silently concatenated.

This cold-start ladder applies to mature \texttt{otr\_ofr\_rv}. It is intentionally
separate from \texttt{BondNewIssue} eligibility.

\paragraph{Mean reversion versus trend.}
The default economic hypothesis is mean reversion after the post-issuance liquidity
premium has stabilized: OFR cheapness relative to OTR is expected to normalize as
benchmark-demand, auction, and initial positioning effects decay.  It is not an
assumption that every 30Y OTR/OFR spread is stationary.  Fresh-issue periods can
trend or gap because of auction supply, dealer inventory, benchmark/index demand,
special repo, and short-squeeze conditions.  Accordingly, the static \texttt{TBondCurve}
family default remains MR only after the pair-specific warm-up; the existing monthly,
point-in-time regime classifier may route a mature eligible pair to the trend engine
when it identifies a trending regime.  An uncertain or cold-start pair is not forced
into either engine and is not traded.

\paragraph{New-issue event strategy (\texttt{BondNewIssue}).}
\texttt{BondNewIssue} is a dedicated event-driven strategy, not a
\texttt{TBondCurve} signal variant. Its instrument is role-based and rebinds across
rolls (new OTR, first OFR, optional second OFR control), so it does not have the
continuous pair history needed for MR/trend z-score entry. Entry is based on a
walk-forward issuance-cohort percentile indexed by OTR age.

The strategy is generalized across CGB and CDB buckets (5Y/10Y/30Y and future
eligible tenors) via \texttt{asset\_class} and \texttt{tenor\_bucket}. Release and
risk gates are evaluated independently per \texttt{(asset\_class, tenor\_bucket)};
approval for one bucket does not unlock any other bucket.

\paragraph{Scope boundary.}
This sub-strategy does not extend the affine curve calibration beyond its validated
1--10Y range.  Mature RV shares \texttt{TBondCurve}/\texttt{CBondCurve} data, risk,
candidate, allocation, and backtest plumbing via \texttt{signal\_variant=otr\_ofr\_rv}.
\texttt{BondNewIssue} has its own loader/backtest path but shares portfolio-level
risk aggregation and execution controls.

\paragraph{Implementation status (2026-08-06).} Config/schema
(\texttt{NewIssueConfig}, \texttt{BondConfig.TBOND\_CURVE\_VARIANTS}), the
point-in-time OTR/OFR universe builder, the \texttt{BondNewIssue-spds.pkl}
artifact, loader/legs wiring, a walk-forward cohort-percentile
\texttt{EventDriven} backtest engine, and the mandatory-gate state machine
(\texttt{DISCOVERED}$\to$\texttt{ORDER\_ELIGIBLE}) are implemented and unit
tested. There is no live order-management/execution venue in this codebase,
so the state machine only certifies up to \texttt{ORDER\_ELIGIBLE}; it does
not submit, fill, or route orders. As documented above, only 30Y CGB has
enough issuance episodes for prototype cohort scoring — all other buckets
remain \texttt{insufficient\_episodes}/\texttt{insufficient\_bucket\_bonds}
by design until independently audited.

\section{Alpha Portfolio $\rightarrow$ Backtest: Style-routed Spread Strategies}

\paragraph{Implementation status (2026-08-03).} The following sections distinguish
the intended methodology from the code currently wired to the dashboard. Individual
backtest ``Auto Monthly'' does classify each spread at the first available observation
of a month using information through that date, then runs an MR or trend engine over a
180-calendar-day warm-up slice ending at that month's end and stitches the in-month
equity segments together. It does \emph{not} preserve a single open trade state across
month boundaries, so it cannot currently implement or audit a true
\texttt{monthly\_style\_change} close. Portfolio Backtest is different again: it loads
the current persisted Alpha allocation, applies today's normalized weights to the full
historical per-asset equity curves, and does not run the monthly style schedule. It is
therefore a current-book historical sanity view, not a causal regime-aware portfolio
backtest. The planned replacement is a stateful, monthly walk-forward simulation with
dated historical allocation snapshots, point-in-time style schedules, next-bar fills,
and explicit rebalance/transaction-cost accounting.

The Alpha Book backtest in \texttt{web/tabs/alpha/} evaluates individual relative-value
spreads and combines the current Portfolio allocation into an indicative book
equity curve. Each spread is assigned exactly one style, mean reversion (MR) or
trend, on the first available trading day of each calendar month. The assignment
is fixed through month-end: the model does not adaptively switch engines within
the month. Candidates are scored separately within their assigned style at that
same review. The resulting frozen score determines eligibility, direction, ranking,
and allocation; it is not an intramonth entry or exit confirmation.

{\small
\setlength{\LTleft}{0pt}
\setlength{\LTright}{0pt}
\setlength{\tabcolsep}{5pt}
\renewcommand{\arraystretch}{1.15}
\begin{longtable}{p{0.16\textwidth} p{0.24\textwidth} p{0.29\textwidth} p{0.25\textwidth}}
\caption{Style-routed spread strategy types.}\label{tab:pairs-strategy-types-overleaf} \\
\toprule
\textbf{Strategy type} & \textbf{Engine} & \textbf{Economic hypothesis} & \textbf{Core entry confirmation} \\
\midrule
\endfirsthead
\toprule
\textbf{Strategy type} & \textbf{Engine} & \textbf{Economic hypothesis} & \textbf{Core entry confirmation} \\
\midrule
\endhead
\midrule
\multicolumn{4}{r}{\textit{Continued on next page}} \\
\midrule
\endfoot
\bottomrule
\endlastfoot
Mean-reverting (MR) & \texttt{engine\_mr.py} & A statistically rich/cheap spread will revert toward its local equilibrium. & A 120-day rolling z-score crosses the entry band. \\
Trend & \texttt{engine\_trend.py} & A persistent directional move should continue once medium-horizon momentum is strong relative to current noise. & A hysteresis crossing of a z-scored momentum signal ($z_t \ge +\theta_z$ for long, $z_t \le -\theta_z$ for short). \\
\end{longtable}
}

\subsection{Regime Selection and Scope}
\label{sec:pairs-regime-selection}

\texttt{curves/calibration/regime.py::compute\_regime\_features} computes four rolling,
60-day indicators on first differences: Kaufman efficiency ratio, a single-scale
R/S Hurst estimate, a variance-ratio proxy, and lag-1 autocorrelation. Each casts a
trend / mean-reversion / neutral vote. A net vote of at least $+2$ selects
\texttt{trending}; at most $-2$ selects \texttt{mean\_reverting}; otherwise the result is
\texttt{uncertain}.

The classifier is invoked only at the monthly review date, using data available on
or before that date. Its result is mapped to canonical \texttt{mr} or \texttt{trend}
and held through month-end. If the result is uncertain or history is insufficient,
the prior monthly style is retained; before the first valid classification, the
spread family's static default applies. If the new style differs from an open
position's style, that position is closed once at the review observation with
\texttt{exit\_reason=monthly\_style\_change}. Regime details are retained in the review
audit trail, not used for daily routing or as a hidden entry filter.

Each monthly record contains the review date, assigned style, classifier label,
vote score/confidence, fallback reason, and the effective trend-signal parameter set
(momentum horizon, robust-volatility window, hysteresis threshold). This makes both
the style and the timing rule reproducible without look-ahead bias.

\subsection{Monthly Candidate Ranking and Allocation}
\label{sec:pairs-monthly-ranking}

Candidate selection and trade timing are deliberately separate decisions. At the
monthly review, MR candidates retain the expected-return/risk score used for
cross-sectional ranking, after stationarity and execution-feasibility checks. Their
daily z-score rule determines whether and when an allocated candidate enters.

Trend candidates use the same combined carry + z-score expected-return score family
as MR candidates; trend does not introduce a separate production scoring formula.
This keeps ranking comparable across regimes while preserving style-specific entry/
exit mechanics in the backtest engine.

An explicit implementation choice remains whether to rank MR and trend candidates in
separate leaderboards or as one pooled cross-sectional list before allocation. The
baseline starts with separate within-regime ranking and validates the pooled
alternative only via family-level walk-forward evidence.

\subsection{Tenor-aware derivative margin proxy}
\label{sec:alpha-tenor-aware-margin}

The portfolio allocator does not have access to a CCP initial-margin file or a
bilateral SIMM engine. A flat percentage of derivative notional is therefore
used only as a fallback for unrecognised leg codes. Recognised IRS legs use a
simple, transparent proxy based on gross leg DV01:
\begin{equation}
 M_{\mathrm{IRS}} = \max\left(
 \frac{\mathrm{DV01}_{\mathrm{gross}}\,s_{\mathrm{tenor}}}{1000},
 \; q_{\min}N_{\mathrm{gross}}
 \right),
\end{equation}
where $M_{\mathrm{IRS}}$ and $N_{\mathrm{gross}}$ are in MM CNY,
$\mathrm{DV01}_{\mathrm{gross}}$ is in thousand CNY per basis point, and
$s_{\mathrm{tenor}}$ is a conservative stress in basis points. The current
prototype buckets use 10, 20, 35, 50, and 75 bp for maximum leg tenors up to
1Y, 3Y, 5Y, 10Y, and longer, respectively, with a 0.25\% gross-notional
liquidity floor. For example, a 6M--1Y spread is charged against the short
end's gross DV01 rather than against the same flat 3\% notional rate used for
all swap spreads.

This is a capital-allocation proxy, not a market quotation. Cleared swaps
should ultimately use the relevant CCP initial-margin model; bilateral swaps
should use the applicable SIMM sensitivities, tenor buckets, correlations,
netting set, liquidity add-ons, and collateral terms. The allocator's
duration-adjusted covariance remains a separate risk-parity input and is not
reused as a margin requirement. Unrecognised derivative legs retain the
legacy flat-rate fallback and should be reviewed before production use.

After ranking, an extreme-score filter removes candidates that imply buying an
already-rich level or selling an already-cheap level relative to assigned direction.
This is the monthly universe-level counterpart of daily chase-prevention bounds.

Correlation, capacity, DV01, and margin constraints are applied after ranking and
the extreme-score filter. The dated allocation snapshot retains the source-data
as-of timestamp, factor values, selected direction, score, and allocation. Neither
the score nor a subsequent change in rank may act as an intramonth entry, exit, or
reversal gate.

\paragraph{Indicator definitions.} Let $s_t$ be the spread level, $\Delta s_t = s_t -
s_{t-1}$ its first difference, and $w=60$ observations (\texttt{DEFAULT\_REGIME\_WINDOW})
the trailing window. All four indicators are computed on the most recent $w$ values of
$\Delta s$ (or, for the efficiency ratio, on the corresponding $w$-step level change):

\begin{equation}
\begin{aligned}
\text{Efficiency ratio: } && ER_t &= \frac{\left|s_t - s_{t-w}\right|}{\displaystyle\sum_{i=t-w+1}^{t} \left|\Delta s_i\right|}, \\[4pt]
\text{Hurst exponent (single-scale R/S): } && Y_i &= \sum_{k=1}^{i}\bigl(\Delta s_k - \overline{\Delta s}\bigr),\quad
R = \max_i Y_i - \min_i Y_i,\quad S = \operatorname{sd}(\Delta s), \\
&& H_t &= \frac{\ln(R/S)}{\ln w}, \\[4pt]
\text{Variance ratio (Lo-MacKinlay): } && k &= \max\!\left(\left\lfloor w/12\right\rfloor,\,2\right),\quad
\hat\mu = \overline{\Delta s},\quad \hat\sigma_a^2 = \frac{1}{w-1}\sum_{i=t-w+1}^{t}(\Delta s_i - \hat\mu)^2, \\
&& \hat\sigma_b^2(k) &= \frac{1}{m}\sum_{i=t-w+k}^{t}\Bigl(\textstyle\sum_{j=0}^{k-1}\Delta s_{i-j} - k\hat\mu\Bigr)^{\!2},\quad
m = k(w-k+1)\Bigl(1-\tfrac{k}{w}\Bigr), \\
&& VR_t &= \hat\sigma_b^2(k)\big/\hat\sigma_a^2, \\[4pt]
\text{Lag-1 autocorrelation: } && AC_t &= \operatorname{corr}\!\left(\Delta s_i,\ \Delta s_{i-1}\right) \text{ over the trailing } w \text{ observations}.
\end{aligned}
\end{equation}

Plain-language reading of each indicator:
\begin{itemize}
\item \textbf{Efficiency ratio $ER_t \in [0,1]$}: ratio of the net $w$-step displacement to
the total path length traveled to get there. Near $1$: the spread moved directly in one
direction with little back-and-forth (efficient/trending). Near $0$: large round-trip
churn relative to the net move (choppy/range-bound, favors mean reversion).
\item \textbf{Hurst exponent $H_t$}: $H_t>0.5$ indicates persistence (an up-move tends to
be followed by another up-move --- trending); $H_t=0.5$ is consistent with a random
walk (no memory); $H_t<0.5$ indicates anti-persistence (moves tend to reverse ---
mean-reverting).
\item \textbf{Variance ratio $VR_t$}: the Lo-MacKinlay (1988) unbiased overlapping
estimator, comparing the per-period variance of $k$-step sums against the 1-step
variance. $VR_t>1$ means variance is growing faster than a random walk would imply
(trending/momentum); $VR_t<1$ means it grows slower (mean-reverting, i.e.\ moves get
walked back). $k$ is kept at $\lfloor w/12\rfloor$ (not $\lfloor w/4\rfloor$) so the
estimator retains enough independent overlapping blocks to stay unbiased at $w=60$.
\item \textbf{Lag-1 autocorrelation $AC_t$}: positive autocorrelation of daily changes
indicates momentum/continuation; negative autocorrelation indicates a reversal
tendency (mean reversion).
\end{itemize}

\paragraph{Voting thresholds.} Each indicator casts one vote in $\{-1, 0, +1\}$
(mean-reverting, neutral, trending). The efficiency ratio and Hurst thresholds are set
\emph{relative to their random-walk null}, not as fixed constants, because both
statistics' expected value under i.i.d.\ noise depends on $w$:
$\mathbb{E}[ER_t \mid \text{i.i.d.}] \approx \sqrt{2/(\pi w)}$ shrinks as $w$ grows, and
$\mathbb{E}[H_t \mid \text{i.i.d.}]$ (the finite-sample R/S bias, calibrated by Monte
Carlo simulation of the same estimator rather than a closed-form correction, since the
classic Anis-Lloyd formula does not track this exact R/S convention closely enough to
be reliable) drifts upward from $0.5$ with $w$. Anchoring the gate to the null keeps a
window change (or a change of instrument, since realised volatility and observation
frequency both shift the effective null) from silently relabeling a flat series:

{\small
\setlength{\LTleft}{0pt}
\setlength{\LTright}{0pt}
\setlength{\tabcolsep}{5pt}
\renewcommand{\arraystretch}{1.15}
\begin{longtable}{p{0.24\textwidth} p{0.22\textwidth} p{0.22\textwidth} p{0.24\textwidth}}
\caption{Regime-vote thresholds per indicator (window-relative; values shown for $w=60$).}\label{tab:pairs-regime-votes-overleaf} \\
\toprule
\textbf{Indicator} & \textbf{Trending vote ($+1$)} & \textbf{Mean-reverting vote ($-1$)} & \textbf{Otherwise} \\
\midrule
\endfirsthead
\toprule
\textbf{Indicator} & \textbf{Trending vote ($+1$)} & \textbf{Mean-reverting vote ($-1$)} & \textbf{Otherwise} \\
\midrule
\endhead
\midrule
\multicolumn{4}{r}{\textit{Continued on next page}} \\
\midrule
\endfoot
\bottomrule
\endlastfoot
Efficiency ratio $ER_t$ & $> 1.6\,\mathbb{E}[ER\mid\mathrm{iid}]$ ($\approx 0.16$ at $w{=}60$) & $< 0.7\,\mathbb{E}[ER\mid\mathrm{iid}]$ ($\approx 0.07$) & neutral ($0$) \\
Hurst exponent $H_t$ & $> \mathbb{E}[H\mid\mathrm{iid}] + 0.05$ ($\approx 0.57$) & $< \mathbb{E}[H\mid\mathrm{iid}] - 0.05$ ($\approx 0.47$) & neutral ($0$) \\
Variance ratio $VR_t$ & $> 1.05$ & $< 0.95$ & neutral ($0$) \\
Lag-1 autocorrelation $AC_t$ & $> 0.05$ & $< -0.05$ & neutral ($0$) \\
\end{longtable}
}

The variance ratio and autocorrelation nulls are window-invariant ($1.0$ and $0.0$
respectively once $VR_t$ uses the bias-corrected estimator above), so their bands stay
fixed constants; only $ER_t$ and $H_t$ use window-relative gates.

The four votes are summed and normalized to a \texttt{regime\_score} $\in [-1, +1]$ (sum
divided by $4$); \texttt{regime\_confidence} is $\left|\texttt{regime\_score}\right|$. The
net (unnormalized) vote sum determines the label: $\ge +2$ votes selects
\texttt{trending}; $\le -2$ votes selects \texttt{mean\_reverting}; any net vote in
$\{-1, 0, +1\}$ is reported as \texttt{uncertain}. This ensemble-voting design means at
least two of the four indicators must agree in the same direction before a definite
regime is assigned; a single strong indicator outvoted by the others yields
\texttt{uncertain} rather than a confident call.

\paragraph{Calibration fix (2026-08-16).} An earlier version of this classifier used
fixed thresholds ($ER>0.35$/$<0.20$, $H>0.55$/$<0.45$) and an uncorrected overlapping
variance-ratio estimator with $k=\lfloor w/4\rfloor$. Simulating a pure random walk
(no genuine trend or mean-reversion signal by construction) through that version showed
it was mislabeled \texttt{mean\_reverting} in roughly 60\% of draws, with a mean net vote
of $-1.25$ out of a possible $[-4,+4]$ range that should average to $0$ on noise. Two of
the four voters were responsible: the efficiency-ratio gate sat above the random-walk
null at $w=60$ (making the trending branch almost unreachable in practice), and the
uncorrected variance-ratio estimator was downward-biased to a mean of $\approx 0.72$
instead of $1.0$ at $k=\lfloor w/4\rfloor=15$. Both defects scaled with $w$, so the bias
was not specific to any one spread or window choice --- it affected every family running
through \texttt{compute\_regime\_features}/\texttt{compute\_regime\_features\_series},
including the monthly style router of \S\ref{sec:pairs-regime-selection} and the hybrid
backtest engine's daily regime gate. The fix above (Lo-MacKinlay variance ratio at
$k=\lfloor w/12\rfloor$, window-relative $ER$/$H$ gates calibrated against each
statistic's random-walk null) reduces the mean vote on simulated random walks to
approximately $-0.2$ and is covered by
\texttt{tests/test\_regime\_hybrid.py::test\_regime\_classifier\_is\_unbiased\_on\_random\_walk}.
Re-running the fixed classifier against several years of CGB-10s30s history now assigns
\texttt{trending} through 2022H2--2023H1 and late 2023 through most of 2024, and returns
to \texttt{trending} from late 2025, consistent with the qualitative trending/mean-reverting
pattern visible in the raw spread chart over that period.

\paragraph{Not a Markov/HMM model.} Both \texttt{compute\_regime\_features} and the
trend-state signal of \S\ref{sec:pairs-trend-model} are deterministic,
rule-based calculations on a rolling window --- there is no hidden-state inference, no
fitted transition matrix, and no probabilistic smoothing. This is a deliberate
``lightweight real-time path'' design choice (see the module docstring). Note that
\texttt{futures/backtest/regime.py} additionally defines a separate
\texttt{RegimeDetector} class built on \texttt{hmmlearn.GaussianHMM} for the futures
strategy stack; the Alpha/pairs regime classifier described here imports only the two
stateless helper functions from that module (\texttt{\_estimate\_hurst} and
\texttt{\_calculate\_variance\_ratio}) and does \textbf{not} use, fit, or depend on that
HMM class.

\subsection{Mean-reversion Model and Defaults}
\label{sec:pairs-mr-model}

For a spread level $s_t$, the MR engine anchors to a 120-observation rolling mean but
scales deviations by a shorter-span EWMA volatility rather than the rolling standard
deviation over the same 120-observation window:

\begin{equation}
z_t = \frac{s_t - \operatorname{mean}_{120}(s)}{\operatorname{ewmstd}_{\mathrm{span}=60}(s)}.
\end{equation}

The engine enters long when $z_t \le -z_{entry}$ and short when
$z_t \ge z_{entry}$. A signal exit is allowed only after \texttt{min\_hold}; it occurs
when the z-score reverts inside the exit band. A $z$-score stop is always active,
including during the minimum holding period. Carry/roll and direction-aware borrow
adjustments accrue in realized PnL and can inform ranking and allocation, but do not
change an MR entry or exit decision.

{\small
\setlength{\LTleft}{0pt}
\setlength{\LTright}{0pt}
\setlength{\tabcolsep}{5pt}
\renewcommand{\arraystretch}{1.15}
\begin{longtable}{p{0.30\textwidth} p{0.16\textwidth} p{0.18\textwidth} p{0.28\textwidth}}
\caption{Mean-reversion strategy defaults.}\label{tab:pairs-mr-defaults-overleaf} \\
\toprule
\textbf{Parameter} & \textbf{Standard default} & \textbf{Tenor-spread preset} & \textbf{Role} \\
\midrule
\endfirsthead
\toprule
\textbf{Parameter} & \textbf{Standard default} & \textbf{Tenor-spread preset} & \textbf{Role} \\
\midrule
\endhead
\midrule
\multicolumn{4}{r}{\textit{Continued on next page}} \\
\midrule
\endfoot
\bottomrule
\endlastfoot
Rolling-mean anchor window & 120 observations & 120 observations & Local equilibrium (fair-value) level; not UI-configurable. \\
EWMA volatility span (\texttt{mr\_vol\_span}) & 60 observations & 60 observations & Dispersion scale for $z_t$; not UI-configurable, but reacts to the current vol regime within its span instead of blending in up to a year of trailing history. \\
\texttt{entry\_z} & 2.0 & 2.5 & Absolute z-score entry threshold. \\
\texttt{exit\_z} & 0.5 & 0.25 & Reversion band used after the minimum hold. \\
\texttt{stop\_z} & 4.0 & 5.0 & Adverse z-score stop. \\
\texttt{min\_hold} & 7 calendar days & 10 calendar days & Minimum holding period for signal exits; stop remains active. \\
\end{longtable}
}

\paragraph{Volatility-scale fix (2026-08-16).} An earlier version scaled $z_t$ by a
plain rolling standard deviation over the same 120-observation window used for the
mean, i.e.\ $z_t = (s_t - \operatorname{mean}_{120}(s))/\operatorname{sd}_{120}(s)$.
Because that window is long enough to span more than one volatility regime, once a
spread's realised volatility fell materially below its trailing-year level, the
rolling std kept the entry threshold anchored to the \emph{old}, higher-vol regime,
so daily moves that were large relative to \emph{current} conditions no longer
produced $|z_t| \ge z_{entry}$ and entries silently stopped even though the spread
remained active in relative terms. Observed on \texttt{LGBCGB-30y}: with the
\texttt{TenorSpread} preset (\texttt{entry\_z}$=2.5$), the rolling-std $z_t$ reached
$|z_t|\ge 2.5$ on 13 days in 2023 but 0 days in 2026, even though the mean 120-day
rolling volatility itself was essentially unchanged between the two years ($\approx$0.028
in both). Replacing the denominator with a 60-span EWMA
volatility, together with lowering \texttt{entry\_z} to $2.0$, raised the full-history
trade count on \texttt{LGBCGB-30y} through \texttt{run\_monthly\_style\_backtest} from
37 to 47 trades and restored entries in 2026 (0 $\to$ 2). The mean anchor is left on
the 120-observation rolling window: an equally short-span EWMA mean was evaluated and
rejected, since a fast-adapting mean chases the level and \emph{reduces} the measured
deviation, which made the effect worse rather than better in exactly the low-vol
periods this fix targets. This applies to both
\texttt{engine\_mr.py::run\_spread\_backtest} (the portfolio-level MR sleeve) and
\texttt{engine\_monthly.py::run\_monthly\_style\_backtest} (the regime-routed
individual-spread and monthly-switch backtest of \S\ref{sec:pairs-regime-selection}),
which previously duplicated the same fixed-window computation independently. This fix
covers backtest \emph{entry timing} only; the separate candidate-\emph{selection}
\texttt{Zscore} computed for the scan/ranking table had the same structural weakness
via a different code path and required its own fix -- see the
``Candidate-selection Zscore volatility fix'' paragraph in
\S\ref{sec:pairs-model-comparison}.

\paragraph{OU-mean consistency fix (2026-08-27).} The volatility-scale fix above left
one inconsistency unaddressed: \texttt{OU\_calibrate} runs an ADF stationarity test on
the raw spread level and, when the series tests stationary, anchors its
candidate-selection \texttt{Zscore} to the AR(1)/OU long-run mean $\theta$ rather than
a sample average -- but \texttt{engine\_mr.py::run\_spread\_backtest} and
\texttt{engine\_monthly.py::run\_monthly\_style\_backtest} computed $z_t$ against a
plain $\operatorname{mean}_{120}(s)$ regardless of that same ADF result, with no
stationarity test of their own. For a genuinely stationary spread this is not merely a
different convention: a rolling mean of a mean-reverting series adds lag and noise
without adding information, since there is no drift for it to track, so the ADF-
qualified OU mean is the statistically appropriate anchor whenever it is available.
This was also the direct cause of the entry-point inconsistency reported against
\texttt{CGB-5s10s} on 2026-08-27, where the live Daily Spread Statistics
\texttt{Zscore} (OU mean, $1.15$) and a chart-level rolling-mean \texttt{Zscore}
disagreed in both magnitude and sign at the same date ($\approx-1$).

Both engines now accept an optional \texttt{ou\_mean} parameter
(\texttt{StatInfo['mean']}, passed by the callback only when
\texttt{StatInfo['stationary']=='YES'}) and blend it with the existing rolling mean via
\texttt{web/tabs/alpha/backtest/\_ou\_mean.py::blended\_mr\_mean}: the OU mean overrides
$\operatorname{mean}_{120}(s)$ for the trailing \texttt{GeneralConfig.STAT\_WINDOW}
(12) months -- the same window \texttt{OU\_calibrate} itself was fit over -- and older
history keeps the plain rolling mean unchanged. A single static OU mean projected
across a spread's \emph{entire} multi-year backtest history was evaluated and
rejected: on \texttt{CGB-5s10s}, whose full history trended well beyond the range the
trailing-12-month OU fit ever saw ($\operatorname{mean}_{120}(s)$ ranged
$0.02$--$0.35$ over the full sample against an OU mean of $0.27$), applying the OU mean
uniformly turned $36$ trades into $643$ by scoring years-old spread levels against
today's fair value. The trailing-window blend keeps that recent-period correction
(the latest $z_t$ moves from $0.45$ under the old rolling-mean anchor to $1.39$, now
consistent with the live snapshot) while leaving older history's trade count
essentially unchanged ($36\to34$). Non-stationary spreads are unaffected: \texttt{ou\_mean}
is only ever passed by the callback when \texttt{stationary=='YES'}, so a trending
series -- where a fixed long-run mean would misread a level shift as an extreme
reading -- keeps the plain rolling-mean anchor exactly as before.

The portfolio-loop backtest (\S\ref{sec:pairs-model-comparison}) receives the same
\texttt{ou\_mean} wiring for its MR-routed assets. Its trend-vs-MR engine choice
remains a static \texttt{style} tag on the persisted allocation snapshot, independent
of \texttt{StatInfo['stationary']} and of the monthly regime classifier in
\S\ref{sec:pairs-regime-selection} -- whether that choice should itself be gated by
stationarity is a separate, larger design question left open for follow-up, not
addressed by this fix.

\subsection{Trend Model and Defaults}
\label{sec:pairs-trend-model}

The production trend signal is a MAD/z-score momentum hysteresis state machine.
\paragraph{Code-symbol compatibility note.} The runtime function is currently named
\texttt{run\_trend\_backtest\_dc} in
\texttt{web/tabs/alpha/backtest/engine\_trend.py} for backward compatibility with
existing callback wiring and tests. Under this methodology, that suffix is legacy:
the intended canonical name is \texttt{run\_trend\_backtest}, with
\texttt{run\_trend\_backtest\_dc} retained only as a temporary alias during
migration.

Let $k$ be the momentum horizon (the engine's \texttt{mom\_window}),
\begin{equation}
m_t = s_t - s_{t-k}.
\end{equation}
Define a robust, continuously updated volatility scale of $m_t$ using a rolling MAD,
\begin{equation}
\sigma_t^{\mathrm{MAD}} = 1.4826\times
\operatorname{median}_{W}\!\left(\left|m_t-
\operatorname{median}_{W}(m)\right|\right),
\end{equation}
and the normalized momentum score
\begin{equation}
z_t = \frac{m_t}{\sigma_t^{\mathrm{MAD}}}.
\end{equation}

Trend state is generated by a Schmitt-trigger hysteresis rule with threshold
$\theta_z>0$:
\begin{equation}
D_t=
\begin{cases}
+1, & z_t \ge +\theta_z, \\
-1, & z_t \le -\theta_z, \\
D_{t-1}, & \text{otherwise},
\end{cases}
\end{equation}
initialized at $D_0=0$ before the first crossing. This keeps state stable between
crossings while avoiding the monthly frozen-threshold staleness of the prior DC
design.

The trade loop is unchanged once $D_t$ is produced: enter long when $D_t>0$, short
when \texttt{allow\_short} and $D_t<0$; exit on opposite-state signal after
\texttt{min\_hold} or immediately on a volatility-scaled trailing stop. Carry and
borrow costs affect realized PnL but do not gate entry.

\paragraph{Deferred research variants.} Two variants remain explicitly research-only:
\emph{trend-following pullback overlay} (delay entry within a confirmed trend to seek
better average fill) and \emph{trend-conditioned fade} (counter-trend turning-point
rule with explicit reversal confirmation). Neither may replace or augment the baseline
without anchored, family-level walk-forward evidence under next-bar execution,
explicit costs, turnover/capacity limits, and a held-out period.

{\small
\setlength{\LTleft}{0pt}
\setlength{\LTright}{0pt}
\setlength{\tabcolsep}{5pt}
\renewcommand{\arraystretch}{1.15}
\begin{longtable}{p{0.24\textwidth} p{0.17\textwidth} p{0.19\textwidth} p{0.32\textwidth}}
\caption{Trend strategy parameters.}\label{tab:pairs-trend-defaults-overleaf} \\
\toprule
\textbf{Parameter} & \textbf{Standard default} & \textbf{Tenor-spread preset} & \textbf{Role} \\
\midrule
\endfirsthead
\toprule
\textbf{Parameter} & \textbf{Standard default} & \textbf{Tenor-spread preset} & \textbf{Role} \\
\midrule
\endhead
\midrule
\multicolumn{4}{r}{\textit{Continued on next page}} \\
\midrule
\endfoot
\bottomrule
\endlastfoot
Momentum horizon \texttt{mom\_window} ($k$) & 20 observations & 20 observations & Lookback for $m_t=s_t-s_{t-k}$. \\
Robust-vol window $W$ & 60 observations & 60 observations & Rolling MAD window for momentum normalization. \\
Hysteresis threshold $\theta_z$ & 1.25 & 1.50 & State flips only when $z_t$ crosses $\pm\theta_z$. \\
\texttt{trailing\_mult} & 1.5 & 2.0 & Multiple of the frozen monthly robust scale. \\
\texttt{min\_hold} & 7 calendar days & 10 calendar days & Minimum hold before an opposite-state signal exit. \\
Re-entry cooldown & 3 bars & 3 bars & Blocks immediate same-direction re-entry after a stop/flip exit. \\
\texttt{allow\_short} & enabled & enabled & Allows short-spread entries. \\
\end{longtable}
}

\subsection{Model Comparison and Deliberate Exclusions}
\label{sec:pairs-model-comparison}

The selected design separates two decisions: monthly style selection and daily
trend-state timing via hysteresis-banded normalized momentum. Table
\ref{tab:pairs-model-comparison-overleaf} records alternatives considered and why they
are not the production entry model.

{
\small
\setlength{\LTleft}{0pt}
\setlength{\LTright}{0pt}
\setlength{\tabcolsep}{5pt}
\renewcommand{\arraystretch}{1.15}
\begin{longtable}{p{0.24\textwidth} p{0.27\textwidth} p{0.38\textwidth}}
\caption{Trend-model comparison and production decision.}\label{tab:pairs-model-comparison-overleaf} \\
\toprule
\textbf{Alternative} & \textbf{Potential benefit} & \textbf{Why it is not the production entry rule} \\
\midrule
\endfirsthead
\toprule
\textbf{Alternative} & \textbf{Potential benefit} & \textbf{Why it is not the production entry rule} \\
\midrule
\endhead
\midrule
\multicolumn{3}{r}{\textit{Continued on next page}} \\
\midrule
\endfoot
\bottomrule
\endlastfoot
Single global momentum threshold without hysteresis & Simple and responsive. & Flips too frequently in noisy intervals; lacks state stability. The baseline keeps a hysteresis band to reduce whipsaw turnover. \\
Monthly frozen threshold on raw spread moves (DC) & Intuitive event logic from local extrema. & Threshold can go stale for weeks after a mid-month volatility shift and requires an extremum-tracking loop. The baseline normalizes momentum continuously with rolling robust volatility and uses one interpretable z-threshold. \\
Trend-following pullback overlay & May improve average entry level inside a confirmed trend. & Adds additional conditions and lag; retained as deferred research, not baseline. \\
Trend-conditioned fade & Can monetize overextended moves. & Counter-trend by construction and can enter before trend exhaustion; retained as deferred research with strict reversal confirmation requirements. \\
Carry-level or candidate-score gate & May favour positive expected carry and highly ranked candidates. & Carry is realized in PnL and candidate scores determine ranking/allocation, but neither establishes a trend reversal. Using either as an entry veto mixes portfolio selection with trade timing and creates a hidden signal filter. \\
Daily adaptive MR/trend routing & Responds quickly to changing market conditions. & Daily switching increases turnover and makes historical attribution and auditability weak. It can switch an open position without a separate trade event. The selected model permits a style change only at the next monthly review. \\
HMM / Markov regime model & Can estimate latent states and transition probabilities. & It adds fitted-state, transition, and retraining risk for limited per-spread histories, while reducing transparency. It is not required for the baseline. HMM research remains separate from this deterministic Alpha model. \\
\end{longtable}
}

The chosen model is therefore not a claim that HMM, pullback, or fade methods cannot
add value. It is a parsimonious baseline with one transparent state process,
continuously normalized momentum, and a complete point-in-time audit trail. Any
alternative must demonstrate incremental out-of-sample net performance, stable
turnover, and comparable interpretability before it replaces this baseline.

For yield-based spread families (\texttt{TenorSpread}, \texttt{SwapSpread},
\texttt{TBondCurve}/\texttt{CBondCurve}, \texttt{TBondSwap}/\texttt{CBondSwap},
\texttt{NetBasis}, \texttt{FuturesSwap}, \texttt{PCASpread}, \texttt{BinarySpread}),
the callback negates the raw spread ($-s_t$) before passing it to either engine, so
that the engine's native $\mathrm{LONG}$ (position $=+1$) already means ``expect the
raw yield-based spread to fall/narrow,'' and native $\mathrm{SHORT}$ already means
``expect the raw spread to rise/widen.'' Prices, $z$-scores, and score series are
sign-restored for display afterwards. This convention must be retained when
reviewing direction, carry, and threshold settings.

\paragraph{Direction-label bug (identified and fixed 2026-07-16).} The individual
backtest callback previously applied an \emph{additional} $\mathrm{LONG}$/$\mathrm{SHORT}$
swap to the trend engine's output only (not to the MR engine's output) after the
sign-restoration step above. Because the engine's direction was already correct
under the $-s_t$ convention, this extra swap double-flipped every trend trade's
displayed direction for yield-based spreads: genuine $\mathrm{SHORT}$ signals (e.g.
during a sustained \texttt{CGB-10s30s} widening move) were shown as $\mathrm{LONG}$
and vice versa, which could be misread as an absence of short signals. PnL figures
were unaffected (they are computed from the underlying position sign, not the
display label). The redundant swap has been removed from
\texttt{web/tabs/alpha/callbacks/backtest\_tab.py}, so trend trades now use the same
(correct) direction convention as MR trades. The trend score-history chart was also
updated to display the same sign convention explicitly.

\paragraph{Candidate \texttt{trend\_state} key bug (identified and fixed 2026-07-16).}
\texttt{curves/refreshers/alpha\_scoring.py} populated the candidate table's
\texttt{trend\_state} column by reading the key \texttt{"state"} from
\texttt{compute\_trend\_signal}'s return dictionary, but that function returns the
value under the key \texttt{"trend\_state"}. The lookup therefore fell back to its
default of \texttt{0.0} for effectively every candidate, silently suppressing the
trend-state signal shown in the candidate scan even when a trend was
confirmed. The key has been corrected.

\paragraph{Candidate-selection Zscore volatility fix (2026-08-16).} The
volatility-scale fix of \S\ref{sec:pairs-mr-model} corrected the backtest engines'
entry-signal scale, but candidate selection computes its own, independent
\texttt{Zscore} for scan/ranking display -- it does not read the backtest engines'
$z_t$ at all. That candidate-side \texttt{Zscore} is \texttt{(spread - mean) / vol},
where \texttt{mean}/\texttt{vol} come from
\texttt{curves/calibration/stat.py::OU\_calibrate} (directly, or via the
\texttt{statAnalysis\_BC}/\texttt{statAnalysis\_BS}/\texttt{statAnalysis\_IRS}
wrappers built on it), computed across roughly ten call sites spanning
\texttt{curves/refreshers/alpha\_snapshot.py}, \texttt{stat.py}, \texttt{irs.py}, and
\texttt{otr\_ofr\_rv.py} for the BondCurve, BondSwap, SwapSpread, TenorSpread,
NetBasis, TermBasis, and FuturesSwap families. \texttt{OU\_calibrate}'s \texttt{vol}
is a \emph{static} standard deviation over the full series passed to it (up to a
252-observation lookback for most families), so it carried the same structural
weakness as the pre-fix backtest engines: once a spread's realised volatility drops
materially below its trailing-year level, its \texttt{Zscore} understates how extreme
the current level actually is, potentially keeping a live candidate below a
\texttt{abs\_zscore}-based ranking/entry bar even though the corrected, EWMA-scaled
backtest $z_t$ for the same instrument on the same day would clear it -- i.e.\ the
present-day snapshot and the historical backtest disagreed about what ``$2\sigma$''
meant for the same spread.

\texttt{OU\_calibrate} now also returns an \texttt{ewm\_vol} column -- an
EWMA(span=$60$) volatility on the same input series, matching
\texttt{MR\_VOL\_SPAN} in \S\ref{sec:pairs-mr-model} -- computed alongside the
existing static \texttt{vol}. The static \texttt{vol} is left unchanged, since it is
also reused for portfolio-level risk sizing (\texttt{web/tabs/alpha/scoring.py},
\texttt{portfolio.py}) where a slower-moving scale is appropriate and redefining it
would have a much larger blast radius than the ranking-consistency issue being fixed
here. Every \texttt{Zscore = (spread - mean) / vol} call site listed above now prefers
\texttt{ewm\_vol}, falling back to the static \texttt{vol} only when \texttt{ewm\_vol}
is unavailable (e.g.\ pre-existing persisted artifacts from before this fix). Verified
on \texttt{LGBCGB-30y}: at the same as-of date, the old static-\texttt{vol} snapshot
\texttt{Zscore} was $-1.79$ against the fixed backtest engine's own $z_t=-1.69$ at that
date; the \texttt{ewm\_vol}-based snapshot \texttt{Zscore} is $-1.73$, materially closer.
The residual gap reflects a separate, intentional difference in the \emph{mean}
anchor -- \texttt{OU\_calibrate} uses the AR(1)/OU-fitted mean when the series tests
stationary, while the backtest engines use a plain rolling(120) mean -- which is not
in scope here. \texttt{curves/refreshers/alpha\_scoring.py}'s regression-residual
\texttt{Zscore} fallback (\texttt{reg\_vol\_resid}, a $30$-day window) was left as-is:
it is only used when the snapshot \texttt{Zscore} is unavailable, and its window is
already short enough not to exhibit the stale-regime bias being corrected.

\subsection{Portfolio Aggregation and Measurement Limitations}

Portfolio mode loads the persisted current Alpha snapshot, normalizes the positive
candidate weights across instruments with available overlapping history, and sums
their weighted daily mark-to-market PnL in basis points. It forwards the MR
entry/exit/stop/minimum-hold controls, but trend assets use the trend engine's
function defaults rather than the individual-screen trend controls. In the current
code path this means the callback calls \texttt{run\_trend\_backtest\_dc}; after
the naming migration above, portfolio and individual paths should call the canonical
\texttt{run\_trend\_backtest} symbol. The resulting
book is a useful \textbf{current-book sanity view}, not a fully parameterized historical
simulation.

The portfolio screen accepts initial capital and a transaction-cost input (default
100 MM and 0.5 bp), but these values are parsed and are not applied to PnL in the
current callback. Likewise, the trend path in portfolio mode does not receive its
spread-type/borrow-cost arguments. Reported portfolio return is consequently
weighted cumulative PnL in bp, rather than capital-normalized net return. The
individual engines report a trade-PnL Sharpe,
$\operatorname{mean}(\mathrm{pnl})/\operatorname{sd}(\mathrm{pnl})\times\sqrt{\min(N_{\mathrm{trades}},20)}$,
whereas the portfolio screen reports a daily-PnL Sharpe annualized by
$\sqrt{252}$. These are not comparable with each other or directly with the
platform-wide return-based convention in \S1.

\paragraph{Regime-aware portfolio target.} The production portfolio backtest must
instead reconstruct the candidate universe, regime assignment, style, and allocation
from dated snapshots available at each monthly review. Each retained instrument carries
its own MR or trend trade state through month boundaries; a changed style or removed
allocation closes the position once at the configured execution price, and a new
position enters only after the applicable engine emits a valid signal. The portfolio
then aggregates daily net returns from contemporaneously active, sized positions and
charges bid--ask, fees, borrow, financing, and rebalance turnover. This design avoids
both retroactive application of current weights and accidental MR routing of stored
``momentum'' style labels.

\paragraph{Trend-conditioned z-score research.} A rule that buys a normalized spread
while its trend state is downward and its moving-average $z$-score remains positive,
and sells while its state is upward and $z$ remains negative, is a counter-trend
turning-point strategy rather than continuation trend following. It can be researched
as a separate, explicitly labeled model, but it should not replace the production
z-momentum continuation baseline by assumption: it may enter before a persistent move
has ended. A safer research sequence compares (i) z-momentum continuation,
(ii) same-direction trend-state continuation with a bounded z-score pullback,
and (iii) the counter-trend rule with an additional reversal confirmation. All variants
must use the economic normalized spread sign, next-bar execution, family-level
walk-forward parameters, transaction costs, and a held-out evaluation period.

\subsection{Validation and Implementation Requirements}

Before increasing risk or selecting family parameters, the implementation must meet
the following validation requirements:

\begin{enumerate}
\item \textbf{Make performance measurement investable first.} Apply instrument- and
direction-specific bid/ask, fees, financing, borrow, and conservative next-bar
execution assumptions on every entry, exit, reversal, and portfolio rebalance;
then make capital and the configured transaction cost operational. Report gross
and net results, turnover, capacity/DV01, and return-based daily Sharpe using one
shared convention.

\item \textbf{Validate with a genuinely walk-forward design.} Freeze the monthly style
and daily trend-signal parameter set (momentum horizon, robust-vol window, hysteresis
threshold, trailing-stop setting) using only information available at each review date;
reserve a final untouched period. Use
purged, embargoed, or anchored walk-forward folds where overlapping holding periods
make ordinary splits optimistic. Compare against always-MR, always-trend, no-trade,
and the rejected alternatives in Table \ref{tab:pairs-model-comparison-overleaf}.

\item \textbf{Calibrate the monthly classifier and trend threshold by family.} The
variance-ratio proxy and single-scale 60-day R/S Hurst estimate have finite-sample
limitations; as of the 2026-08-16 fix described in
\S\ref{sec:pairs-regime-selection}, the classifier corrects for the resulting
random-walk bias generically (Lo-MacKinlay variance ratio, window-relative $ER$/$H$
gates), but the underlying single-scale R/S Hurst estimator and a $60$-day window's
limited power to resolve multi-month regime swings remain approximations. Calibrate
classifier votes and the trend hysteresis threshold ($\theta_z$) to instrument-family
null distributions and walk-forward results, and consider a longer window (120--250
days) for families where multi-month regimes dominate. The review output must always
report confidence, fallback reason, robust scale summary, and effective threshold
configuration.

\item \textbf{Enforce one causal trend-state implementation.} All trend paths must
call the same hysteresis-banded z-momentum signal generator. Test positive,
inverted, near-zero, and zero-crossing spreads; verify no entry occurs before the
first hysteresis crossing; and verify that carry and candidate score cannot gate an
entry or an opposite-state exit.

\item \textbf{Estimate parameters by spread family, not individual in-sample fit.}
Estimate an out-of-sample OU half-life/stationarity diagnostic per instrument; trade
MR only when its half-life fits the intended 1-week to 1-month holding horizon. For
trend, select momentum horizon, robust-volatility lookback, hysteresis threshold,
and trailing-stop multiple by conservative family-level grids, scoring stability
across folds rather than the single best historical Sharpe.

\item \textbf{Improve portfolio construction after signals are credible.} Preserve the
historically available allocation at each rebalance; volatility-target gross DV01;
cap key-rate/issuer/leg and correlated-family exposure; net shared legs; and apply
drawdown, liquidity, and event-risk overlays. Add cross-sectional rich/cheap
signals within a spread family to reduce common level-factor risk that a collection
of independent time-series signals cannot diversify away.
\end{enumerate}
